% ============================================================================
% CAPÍTULO 7: PROTEIN SYSTEMS
% Bioinformática para Biologia de Sistemas
% ============================================================================

% ============================================================================
% TEMPLATE BASE PARA APRESENTAÇÕES EM BEAMER
% Bioinformática para Biologia de Sistemas
% ============================================================================

\documentclass[12pt, aspectratio=169]{beamer}

% ============================================================================
% PACOTES ESSENCIAIS
% ============================================================================
\usepackage[utf8]{inputenc}
\usepackage[T1]{fontenc}
\usepackage[brazilian]{babel}
\usepackage{amsmath, amssymb, amsthm}
\usepackage{graphicx}
\usepackage{xcolor}
\usepackage{tikz}
\usetikzlibrary{arrows.meta, shapes, positioning, calc, decorations.pathreplacing, decorations.shapes, decorations.pathmorphing}
\usepackage{listings}
\usepackage{booktabs}
\usepackage{multirow}
\usepackage{hyperref}
\usepackage{chemfig}
\usepackage{chemarr}

% ============================================================================
% CONFIGURAÇÃO DO TEMA
% ============================================================================
\usetheme{Madrid}
\usecolortheme{default}

% Cores personalizadas
\definecolor{azulescuro}{RGB}{0, 51, 102}
\definecolor{azulclaro}{RGB}{51, 153, 255}
\definecolor{verdeacido}{RGB}{102, 204, 0}
\definecolor{laranjacel}{RGB}{255, 153, 51}
\definecolor{roxodna}{RGB}{153, 51, 255}

\setbeamercolor{structure}{fg=azulescuro}
\setbeamercolor{palette primary}{bg=azulescuro,fg=white}
\setbeamercolor{palette secondary}{bg=azulclaro,fg=white}
\setbeamercolor{palette tertiary}{bg=azulescuro,fg=white}
\setbeamercolor{palette quaternary}{bg=azulclaro,fg=white}
\setbeamercolor{block title}{bg=azulescuro,fg=white}
\setbeamercolor{block body}{bg=azulclaro!10,fg=black}
\setbeamercolor{block title alerted}{bg=red!80,fg=white}
\setbeamercolor{block body alerted}{bg=red!10,fg=black}
\setbeamercolor{block title example}{bg=verdeacido!80,fg=white}
\setbeamercolor{block body example}{bg=verdeacido!10,fg=black}

% ============================================================================
% CONFIGURAÇÃO DE CÓDIGO
% ============================================================================
\lstset{
    language=Python,
    basicstyle=\ttfamily\small,
    keywordstyle=\color{azulescuro}\bfseries,
    commentstyle=\color{verdeacido},
    stringstyle=\color{laranjacel},
    numbers=left,
    numberstyle=\tiny\color{gray},
    stepnumber=1,
    numbersep=5pt,
    backgroundcolor=\color{white},
    showspaces=false,
    showstringspaces=false,
    showtabs=false,
    frame=single,
    rulecolor=\color{black},
    tabsize=2,
    captionpos=b,
    breaklines=true,
    breakatwhitespace=false,
    escapeinside={\%*}{*)},
    xleftmargin=10pt,
    xrightmargin=10pt
}

% Definir estilo para R
\lstdefinestyle{Rstyle}{
    language=R,
    basicstyle=\ttfamily\small,
    keywordstyle=\color{azulescuro}\bfseries,
    commentstyle=\color{verdeacido},
    stringstyle=\color{laranjacel}
}

% Definir estilo para MATLAB
\lstdefinestyle{MATLABstyle}{
    language=Matlab,
    basicstyle=\ttfamily\small,
    keywordstyle=\color{azulescuro}\bfseries,
    commentstyle=\color{verdeacido},
    stringstyle=\color{laranjacel}
}

% ============================================================================
% COMANDOS PERSONALIZADOS
% ============================================================================

% Comando para equações destacadas
\newcommand{\destaque}[1]{%
    \begin{alertblock}{}
        \begin{center}
            #1
        \end{center}
    \end{alertblock}
}

% Comando para definições
\newcommand{\definicao}[2]{%
    \begin{block}{Definição: #1}
        #2
    \end{block}
}

% Comando para exemplos
\newcommand{\exemplo}[2]{%
    \begin{exampleblock}{Exemplo: #1}
        #2
    \end{exampleblock}
}

% Comando para observações importantes
\newcommand{\importante}[1]{%
    \begin{alertblock}{Importante!}
        #1
    \end{alertblock}
}

% Comandos matemáticos úteis
\newcommand{\R}{\mathbb{R}}
\newcommand{\N}{\mathbb{N}}
\newcommand{\Z}{\mathbb{Z}}
\newcommand{\dif}{\mathrm{d}}
\newcommand{\parcial}[2]{\frac{\partial #1}{\partial #2}}
\newcommand{\derivada}[2]{\frac{\dif #1}{\dif #2}}
\newcommand{\vect}[1]{\boldsymbol{#1}}

% Comandos biológicos
\newcommand{\gene}[1]{\textit{#1}}
\newcommand{\proteina}[1]{\textsc{#1}}
\newcommand{\especie}[1]{\textit{#1}}

% ============================================================================
% CONFIGURAÇÃO DE RODAPÉ
% ============================================================================
\setbeamertemplate{footline}{
  \leavevmode%
  \hbox{%
  \begin{beamercolorbox}[wd=.33\paperwidth,ht=2.25ex,dp=1ex,center]{author in head/foot}%
    \usebeamerfont{author in head/foot}\insertshortauthor
  \end{beamercolorbox}%
  \begin{beamercolorbox}[wd=.34\paperwidth,ht=2.25ex,dp=1ex,center]{title in head/foot}%
    \usebeamerfont{title in head/foot}\insertshorttitle
  \end{beamercolorbox}%
  \begin{beamercolorbox}[wd=.33\paperwidth,ht=2.25ex,dp=1ex,right]{date in head/foot}%
    \usebeamerfont{date in head/foot}\insertshortdate{}\hspace*{2em}
    \insertframenumber{} / \inserttotalframenumber\hspace*{2ex}
  \end{beamercolorbox}}%
  \vskip0pt%
}

% ============================================================================
% CONFIGURAÇÃO DE NAVEGAÇÃO
% ============================================================================
\setbeamertemplate{navigation symbols}{}

% ============================================================================
% CONFIGURAÇÃO DE ITEMIZE/ENUMERATE
% ============================================================================
\setbeamertemplate{itemize items}[circle]
\setbeamertemplate{enumerate items}[default]

% ============================================================================
% FIM DO TEMPLATE
% ============================================================================


% Bibliotecas adicionais
\usetikzlibrary{decorations.pathmorphing}

% ============================================================================
% INFORMAÇÕES DO DOCUMENTO
% ============================================================================
\title[Cap. 7: Protein Systems]{Capítulo 7: Protein Systems}
\subtitle{Bioinformática para Biologia de Sistemas}
\author{Ney Lemke}
\institute{DBF-Unesp}
\date{\today}

\begin{document}

% ============================================================================
% SLIDE DE TÍTULO
% ============================================================================
\begin{frame}
\titlepage
\end{frame}

% ============================================================================
% SUMÁRIO
% ============================================================================
\begin{frame}{Sumário}
\tableofcontents
\end{frame}

% ============================================================================
% SEÇÃO 1: INTRODUÇÃO
% ============================================================================
\section{Introdução}

\begin{frame}{Visão Geral do Capítulo}
\begin{block}{Objetivos de Aprendizagem}
\begin{itemize}
    \item Aprofundar conhecimento sobre estrutura de proteínas
    \item Dominar cinética enzimática avançada
    \item Compreender cascatas de sinalização celular
    \item Estudar modificações pós-traducionais
    \item Explorar métodos de proteômica
    \item Aplicar ferramentas bioinformáticas para análise de proteínas
\end{itemize}
\end{block}

\vspace{0.3cm}

\importante{
Proteínas são as máquinas moleculares da vida: executam, regulam e coordenam praticamente todos os processos celulares
}
\end{frame}

\begin{frame}{Importância das Proteínas em Sistemas Biológicos}
\begin{columns}
\column{0.5\textwidth}
\textbf{Funções Estruturais:}
\begin{itemize}
    \item Colágeno (tecido conectivo)
    \item Queratina (cabelo, unhas)
    \item Actina/Miosina (músculo)
    \item Histonas (empacotamento DNA)
\end{itemize}

\vspace{0.3cm}

\textbf{Funções Catalíticas:}
\begin{itemize}
    \item Enzimas metabólicas
    \item DNA/RNA polimerases
    \item Proteases
    \item Quinases
\end{itemize}

\column{0.5\textwidth}
\textbf{Funções Regulatórias:}
\begin{itemize}
    \item Fatores de transcrição
    \item Hormônios peptídicos
    \item Receptores de membrana
    \item Proteínas G
\end{itemize}

\vspace{0.3cm}

\textbf{Funções de Sinalização:}
\begin{itemize}
    \item Anticorpos
    \item Citocinas
    \item Neurotransmissores
    \item Mensageiros intracelulares
\end{itemize}
\end{columns}
\end{frame}

\begin{frame}{Papel Central das Proteínas}
\begin{center}
\begin{tikzpicture}[node distance=2.5cm, scale=0.85, transform shape]
    % Central node
    \node[draw, circle, fill=verdeacido!50, text width=2cm, align=center, font=\large\bfseries] (prot) {PROTEÍNAS};

    % Surrounding functions
    \node[draw, ellipse, fill=azulclaro!30, above left=1.5cm and 2cm of prot, text width=2cm, align=center] (est) {Estrutura\\Celular};
    \node[draw, ellipse, fill=laranjacel!30, above right=1.5cm and 2cm of prot, text width=2cm, align=center] (cat) {Catálise\\Enzimática};
    \node[draw, ellipse, fill=roxodna!30, below right=1.5cm and 2cm of prot, text width=2cm, align=center] (reg) {Regulação\\Gênica};
    \node[draw, ellipse, fill=yellow!50, below left=1.5cm and 2cm of prot, text width=2cm, align=center] (sin) {Sinalização\\Celular};

    % Arrows
    \draw[->, thick, azulescuro] (prot) -- (est);
    \draw[->, thick, azulescuro] (prot) -- (cat);
    \draw[->, thick, azulescuro] (prot) -- (reg);
    \draw[->, thick, azulescuro] (prot) -- (sin);
\end{tikzpicture}
\end{center}

\destaque{
Proteínas são o elo entre genótipo (DNA) e fenótipo (função)
}
\end{frame}

% ============================================================================
% SEÇÃO 2: ESTRUTURA DE PROTEÍNAS
% ============================================================================
\section{Estrutura de Proteínas}

\begin{frame}{7.1 Estrutura de Proteínas - Revisão}
\begin{block}{Quatro Níveis Estruturais}
\begin{enumerate}
    \item \textbf{Primária:} sequência linear de aminoácidos
    \item \textbf{Secundária:} arranjos locais ($\alpha$-hélice, $\beta$-folha)
    \item \textbf{Terciária:} estrutura 3D completa de uma cadeia
    \item \textbf{Quaternária:} arranjo de múltiplas subunidades
\end{enumerate}
\end{block}

\vspace{0.3cm}

\begin{center}
\begin{tikzpicture}[scale=0.7]
    % Primary
    \draw[thick, azulescuro] (0,0) -- (2,0);
    \foreach \x in {0.4, 0.8, 1.2, 1.6}
        \fill[verdeacido] (\x,0) circle (2pt);
    \node[below, font=\tiny] at (1,-0.3) {1ª};

    % Secondary
    \begin{scope}[xshift=3cm]
        \draw[thick, azulescuro, decorate, decoration={snake, segment length=3mm, amplitude=2mm}] (0,0) -- (2,0);
        \node[below, font=\tiny] at (1,-0.5) {2ª};
    \end{scope}

    % Tertiary
    \begin{scope}[xshift=6cm]
        \draw[thick, azulescuro, rounded corners=5pt] (0,0) -- (0.5,0.3) -- (1,0.1) -- (1.5,0.4) -- (2,0);
        \node[below, font=\tiny] at (1,-0.3) {3ª};
    \end{scope}

    % Quaternary
    \begin{scope}[xshift=9cm]
        \draw[thick, verdeacido, fill=verdeacido!20, rounded corners=3pt] (0,0) rectangle (0.8,0.6);
        \draw[thick, laranjacel, fill=laranjacel!20, rounded corners=3pt] (0.9,0) rectangle (1.7,0.6);
        \node[below, font=\tiny] at (0.85,-0.3) {4ª};
    \end{scope}
\end{tikzpicture}
\end{center}
\end{frame}

\begin{frame}{Termodinâmica do Dobramento Proteico}
\begin{block}{Energia Livre de Gibbs}
O dobramento proteico é governado pela mudança de energia livre:
\begin{equation*}
\Delta G = \Delta H - T\Delta S
\end{equation*}

Onde:
\begin{itemize}
    \item $\Delta G < 0$: processo espontâneo (proteína dobrada é estável)
    \item $\Delta H$: entalpia (ligações, interações)
    \item $T\Delta S$: entropia (desordem)
\end{itemize}
\end{block}

\vspace{0.3cm}

\begin{columns}
\column{0.5\textwidth}
\textbf{Contribuições para $\Delta H$:}
\begin{itemize}
    \item Ligações de hidrogênio
    \item Interações eletrostáticas
    \item Van der Waals
    \item Pontes dissulfeto
\end{itemize}

\column{0.5\textwidth}
\textbf{Efeito Hidrofóbico:}
\begin{itemize}
    \item Resíduos hidrofóbicos no interior
    \item Resíduos hidrofílicos na superfície
    \item Minimiza contato água-hidrofóbicos
    \item Principal força motriz
\end{itemize}
\end{columns}
\end{frame}

\begin{frame}{Diagrama de Ramachandran}
\begin{columns}
\column{0.5\textwidth}
\definicao{Ramachandran Plot}{
Gráfico que mostra as combinações permitidas dos ângulos diedros $\phi$ (phi) e $\psi$ (psi) da cadeia principal proteica.
}

\vspace{0.3cm}

\textbf{Regiões Favoráveis:}
\begin{itemize}
    \item \textbf{$\alpha$-hélice:}\\
    $\phi \approx -60°$, $\psi \approx -45°$
    \item \textbf{$\beta$-folha:}\\
    $\phi \approx -120°$, $\psi \approx +120°$
    \item \textbf{Hélice esquerda:}\\
    $\phi \approx +60°$, $\psi \approx +45°$
\end{itemize}

\column{0.5\textwidth}
\begin{center}
\begin{tikzpicture}[scale=0.7]
    % Axes
    \draw[->, thick] (-3,0) -- (3,0) node[right] {$\phi$};
    \draw[->, thick] (0,-3) -- (0,3) node[above] {$\psi$};

    % Grid
    \foreach \x in {-180,-90,0,90,180}
        \draw (\x/60,-0.1) -- (\x/60,0.1);
    \foreach \y in {-180,-90,0,90,180}
        \draw (-0.1,\y/60) -- (0.1,\y/60);

    % Alpha helix region
    \fill[azulclaro!50, opacity=0.7] (-1.5,-1) ellipse (0.8 and 0.6);
    \node[font=\tiny] at (-1.5,-1) {$\alpha$};

    % Beta sheet region
    \fill[verdeacido!50, opacity=0.7] (-2,2) ellipse (0.6 and 0.8);
    \node[font=\tiny] at (-2,2) {$\beta$};

    % Left-handed helix (rare)
    \fill[laranjacel!50, opacity=0.7] (1.5,1) ellipse (0.5 and 0.5);
    \node[font=\tiny] at (1.5,1) {L};

    % Axis labels
    \node[below, font=\tiny] at (-3,-0.2) {-180°};
    \node[below, font=\tiny] at (0,-0.2) {0°};
    \node[below, font=\tiny] at (3,-0.2) {180°};
    \node[left, font=\tiny] at (-0.2,-3) {-180°};
    \node[left, font=\tiny] at (-0.2,0) {0°};
    \node[left, font=\tiny] at (-0.2,3) {180°};
\end{tikzpicture}
\end{center}
\end{columns}

\importante{
Glicina (sem cadeia lateral) permite todos os ângulos. Prolina (imino ácido) restringe $\phi \approx -60°$
}
\end{frame}

\begin{frame}{Domínios e Motivos Proteicos}
\begin{columns}
\column{0.5\textwidth}
\begin{block}{Domínios}
Unidades estruturais e funcionais independentes:
\begin{itemize}
    \item Dobram-se independentemente
    \item 50-200 aminoácidos
    \item Evolução modular
    \item Exemplos: domínio SH2, domínio quinase
\end{itemize}
\end{block}

\begin{exampleblock}{Motivos}
Padrões estruturais recorrentes:
\begin{itemize}
    \item Hélice-volta-hélice
    \item Dedos de zinco
    \item Barril TIM
    \item Leucine zipper
\end{itemize}
\end{exampleblock}

\column{0.5\textwidth}
\begin{center}
\begin{tikzpicture}[scale=0.65]
    % Protein with multiple domains
    \draw[thick, fill=azulclaro!30, rounded corners=5pt] (0,0) rectangle (2,1.5);
    \node at (1,0.75) {\tiny Domínio 1};

    \draw[thick, fill=verdeacido!30, rounded corners=5pt] (2.2,0) rectangle (4.2,1.5);
    \node at (3.2,0.75) {\tiny Domínio 2};

    \draw[thick, fill=laranjacel!30, rounded corners=5pt] (4.4,0) rectangle (6,1.5);
    \node at (5.2,0.75) {\tiny Domínio 3};

    % Linker regions
    \draw[thick, azulescuro] (2,0.75) -- (2.2,0.75);
    \draw[thick, azulescuro] (4.2,0.75) -- (4.4,0.75);

    \node[below, font=\tiny] at (3,-0.3) {Proteína Multidomínio};

    % Motif example
    \begin{scope}[yshift=-3cm]
        % Helix-turn-helix
        \draw[thick, verdeacido, decorate, decoration={snake, segment length=2mm, amplitude=1.5mm}] (0,0) -- (1.5,0);
        \draw[thick, azulescuro] (1.5,0) to [out=45, in=180] (2,0.5);
        \draw[thick, laranjacel, decorate, decoration={snake, segment length=2mm, amplitude=1.5mm}] (2,0.5) -- (3.5,0.5);
        \node[below, font=\tiny] at (1.75,-0.5) {Motivo Hélice-Volta-Hélice};
    \end{scope}
\end{tikzpicture}
\end{center}
\end{columns}
\end{frame}

\begin{frame}{Proteínas Intrinsecamente Desordenadas (IDPs)}
\begin{block}{Características}
Proteínas ou regiões sem estrutura 3D definida em condições fisiológicas:
\begin{itemize}
    \item Alto conteúdo de resíduos polares e carregados
    \item Baixo conteúdo hidrofóbico
    \item Flexibilidade conformacional
    \item 30-40\% do proteoma eucarioto
\end{itemize}
\end{block}

\vspace{0.3cm}

\begin{columns}
\column{0.5\textwidth}
\textbf{Funções:}
\begin{itemize}
    \item Sinalização celular
    \item Regulação da transcrição
    \item Montagem de complexos
    \item Promiscuidade de ligação
\end{itemize}

\column{0.5\textwidth}
\begin{exampleblock}{Exemplos}
\begin{itemize}
    \item p53 (supressor tumoral)
    \item $\alpha$-sinucleína (Parkinson)
    \item Tau (Alzheimer)
    \item Histonas (caudas N-terminais)
\end{itemize}
\end{exampleblock}
\end{columns}

\vspace{0.3cm}

\importante{
IDPs desafiam o paradigma estrutura-função: função sem estrutura fixa!
}
\end{frame}

\begin{frame}{Protein Data Bank (PDB)}
\begin{columns}
\column{0.5\textwidth}
\begin{block}{O que é o PDB?}
\begin{itemize}
    \item Repositório global de estruturas 3D
    \item URL: \url{https://www.rcsb.org}
    \item Mais de 200,000 estruturas
    \item Livre acesso
    \item Formatos: PDB, mmCIF
\end{itemize}
\end{block}

\textbf{Técnicas Experimentais:}
\begin{itemize}
    \item Cristalografia de raios-X (85\%)
    \item Ressonância magnética nuclear (NMR)
    \item Microscopia crioeletrônica (Cryo-EM)
    \item Modelagem computacional
\end{itemize}

\column{0.5\textwidth}
\begin{exampleblock}{Informações no PDB}
\begin{itemize}
    \item Coordenadas atômicas
    \item Estrutura secundária
    \item Ligantes e cofatores
    \item Metadados experimentais
    \item Literatura associada
\end{itemize}
\end{exampleblock}

\vspace{0.3cm}

\begin{alertblock}{Resolução}
Medida em Ångströms (\AA):
\begin{itemize}
    \item Alta resolução: $<$ 2 \AA
    \item Média resolução: 2-3 \AA
    \item Baixa resolução: $>$ 3 \AA
\end{itemize}
\end{alertblock}
\end{columns}
\end{frame}

\begin{frame}[fragile]{Análise de Estruturas com BioPython}
\begin{lstlisting}[language=Python, caption=Parser de arquivos PDB, basicstyle=\ttfamily\footnotesize]
from Bio.PDB import PDBParser, PDBIO, Select
from Bio.PDB.DSSP import DSSP
import numpy as np

# Carregar estrutura
parser = PDBParser(QUIET=True)
structure = parser.get_structure("1UBQ", "1ubq.pdb")

# Extrair informacoes basicas
model = structure[0]
for chain in model:
    print(f"Cadeia {chain.id}: {len(chain)} residuos")

    # Calcular centro de massa
    coords = []
    for residue in chain:
        if residue.id[0] == ' ':  # residuos padrao
            for atom in residue:
                coords.append(atom.coord)

    coords = np.array(coords)
    center_mass = coords.mean(axis=0)
    print(f"Centro de massa: {center_mass}")

# Analisar estrutura secundaria (requer DSSP)
dssp = DSSP(model, "1ubq.pdb")
secondary = [dssp[key][2] for key in dssp.keys()]
print(f"Helices: {secondary.count('H')}")
print(f"Folhas beta: {secondary.count('E')}")
\end{lstlisting}
\end{frame}

\begin{frame}[fragile]{Visualização de Proteínas com PyMOL}
\begin{lstlisting}[language=Python, caption=Script PyMOL para visualização, basicstyle=\ttfamily\footnotesize]
from pymol import cmd

# Carregar estrutura
cmd.load("1ubq.pdb")

# Colorir por estrutura secundaria
cmd.color("red", "ss h")      # helices em vermelho
cmd.color("yellow", "ss s")   # folhas em amarelo
cmd.color("green", "ss l+''") # loops em verde

# Mostrar superficie
cmd.show("surface")
cmd.set("transparency", 0.5)

# Destacar residuos ativos (exemplo: His68)
cmd.select("site", "resi 68")
cmd.show("sticks", "site")
cmd.color("magenta", "site")

# Ajustar visualizacao
cmd.bg_color("white")
cmd.set("ray_shadows", 0)

# Salvar imagem
cmd.png("protein_structure.png", width=1200, dpi=300, ray=1)

print("Visualizacao criada!")
\end{lstlisting}
\end{frame}

% ============================================================================
% SEÇÃO 3: CINÉTICA ENZIMÁTICA AVANÇADA
% ============================================================================
\section{Cinética Enzimática Avançada}

\begin{frame}{7.2 Cinética Enzimática - Revisão de Michaelis-Menten}
\begin{block}{Equação de Michaelis-Menten}
\begin{equation*}
v = \frac{V_{max}[S]}{K_M + [S]}
\end{equation*}

Derivada do mecanismo:
\begin{equation*}
E + S \underset{k_{-1}}{\overset{k_1}{\rightleftharpoons}} ES \overset{k_{cat}}{\longrightarrow} E + P
\end{equation*}
\end{block}

\vspace{0.3cm}

\begin{columns}
\column{0.5\textwidth}
\textbf{Parâmetros:}
\begin{itemize}
    \item $K_M = \frac{k_{-1} + k_{cat}}{k_1}$
    \item $V_{max} = k_{cat}[E]_0$
    \item $k_{cat}$: número de turnover
\end{itemize}

\column{0.5\textwidth}
\textbf{Eficiência Catalítica:}
\begin{equation*}
\frac{k_{cat}}{K_M} \quad \text{(M}^{-1}\text{s}^{-1}\text{)}
\end{equation*}

\begin{itemize}
    \item Limite de difusão: $10^8$-$10^9$ M$^{-1}$s$^{-1}$
    \item Enzimas perfeitas: anidrase carbônica, catalase
\end{itemize}
\end{columns}
\end{frame}

\begin{frame}{Linearização: Gráfico de Lineweaver-Burk}
\begin{columns}
\column{0.5\textwidth}
\begin{block}{Duplo Recíproco}
Invertendo a equação de MM:
\begin{equation*}
\frac{1}{v} = \frac{K_M}{V_{max}} \cdot \frac{1}{[S]} + \frac{1}{V_{max}}
\end{equation*}

Forma linear: $y = mx + b$
\begin{itemize}
    \item Intercepto-y: $1/V_{max}$
    \item Intercepto-x: $-1/K_M$
    \item Inclinação: $K_M/V_{max}$
\end{itemize}
\end{block}

\column{0.5\textwidth}
\begin{center}
\begin{tikzpicture}[scale=0.7]
    % Axes
    \draw[->, thick] (-2,0) -- (4,0) node[right, font=\small] {$1/[S]$};
    \draw[->, thick] (0,-0.5) -- (0,4) node[above, font=\small] {$1/v$};

    % Line
    \draw[thick, azulescuro, domain=-1.5:3.5] plot (\x, {0.5*\x + 1.5});

    % Intercepts
    \fill[red] (0,1.5) circle (2pt);
    \node[left, font=\tiny] at (-0.2,1.5) {$1/V_{max}$};

    \fill[red] (-3,0) circle (2pt);
    \node[below, font=\tiny] at (-3,-0.2) {$-1/K_M$};

    % Slope annotation
    \draw[dashed, verdeacido] (1,2) -- (2,2.5);
    \node[right, font=\tiny] at (2,2.5) {Inclinação = $K_M/V_{max}$};
\end{tikzpicture}
\end{center}

\begin{alertblock}{Desvantagens}
\begin{itemize}
    \item Amplifica erros em $v$ baixo
    \item Distorce distribuição de erros
\end{itemize}
\end{alertblock}
\end{columns}
\end{frame}

\begin{frame}{Gráfico de Eadie-Hofstee}
\begin{columns}
\column{0.5\textwidth}
\begin{block}{Transformação Alternativa}
Dividindo MM por $[S]$:
\begin{equation*}
v = -K_M \cdot \frac{v}{[S]} + V_{max}
\end{equation*}

Vantagens:
\begin{itemize}
    \item Melhor distribuição de erros
    \item Detecta desvios de MM
    \item Identifica cooperatividade
\end{itemize}
\end{block}

\column{0.5\textwidth}
\begin{center}
\begin{tikzpicture}[scale=0.7]
    % Axes
    \draw[->, thick] (0,0) -- (4,0) node[right, font=\small] {$v/[S]$};
    \draw[->, thick] (0,0) -- (0,4) node[above, font=\small] {$v$};

    % Line (negative slope)
    \draw[thick, azulescuro, domain=0:3.5] plot (\x, {3.5 - 0.8*\x});

    % Intercepts
    \fill[red] (0,3.5) circle (2pt);
    \node[left, font=\tiny] at (-0.2,3.5) {$V_{max}$};

    \fill[red] (4.375,0) circle (2pt);
    \node[below, font=\tiny] at (4.375,-0.2) {$V_{max}/K_M$};

    % Slope
    \node[right, font=\tiny, text width=2cm] at (2,2) {Inclinação = $-K_M$};
\end{tikzpicture}
\end{center}
\end{columns}

\vspace{0.3cm}

\begin{exampleblock}{Comparação de Métodos}
Hoje em dia, ajuste não-linear direto (regressão) é preferível às linearizações
\end{exampleblock}
\end{frame}

\begin{frame}{Ligação Cooperativa e Equação de Hill}
\begin{block}{Cooperatividade em Enzimas Oligoméricas}
Quando ligação de substrato afeta afinidade de outros sítios:
\begin{equation*}
v = \frac{V_{max}[S]^n}{K_d^n + [S]^n}
\end{equation*}

Onde $n$ é o coeficiente de Hill:
\begin{itemize}
    \item $n = 1$: sem cooperatividade (Michaeliano)
    \item $n > 1$: cooperatividade positiva
    \item $n < 1$: cooperatividade negativa
\end{itemize}
\end{block}

\exemplo{Hemoglobina}{
Ligação de O$_2$ é cooperativa: $n \approx 2.8$\\
Primeira molécula de O$_2$ facilita ligação das demais (efeito sigmoide)
}
\end{frame}

\begin{frame}{Gráfico de Hill}
\begin{columns}
\column{0.5\textwidth}
\begin{center}
\begin{tikzpicture}[scale=0.7]
    % Axes
    \draw[->, thick] (0,0) -- (5,0) node[right, font=\small] {$[S]$};
    \draw[->, thick] (0,0) -- (0,4) node[above, font=\small] {$v$};

    % Curves for different n
    \draw[thick, azulescuro, domain=0:4.8, samples=100] plot (\x, {3.5*\x/(1+\x)});
    \node[right, font=\tiny] at (3,2.5) {$n=1$};

    \draw[thick, verdeacido, domain=0:4.8, samples=100] plot (\x, {3.5*\x^2/(1+\x^2)});
    \node[right, font=\tiny] at (2,1.5) {$n=2$};

    \draw[thick, laranjacel, domain=0:4.8, samples=100] plot (\x, {3.5*\x^4/(1+\x^4)});
    \node[right, font=\tiny] at (1.3,1) {$n=4$};

    % Vmax line
    \draw[dashed, red] (0,3.5) -- (5,3.5);
    \node[left, font=\tiny] at (0,3.5) {$V_{max}$};
\end{tikzpicture}
\end{center}

\column{0.5\textwidth}
\begin{block}{Gráfico de Hill Linearizado}
\begin{equation*}
\log\frac{v}{V_{max}-v} = n\log[S] - \log K_d
\end{equation*}

Inclinação = $n$ (coeficiente de Hill)
\end{block}

\vspace{0.3cm}

\importante{
$n > 1$ resulta em resposta switch-like: pequena mudança em $[S]$ causa grande mudança em $v$
}
\end{columns}
\end{frame}

\begin{frame}{Regulação Alostérica: Modelos MWC e KNF}
\begin{columns}
\column{0.5\textwidth}
\begin{block}{Modelo MWC}
\textbf{Monod-Wyman-Changeux}
\begin{itemize}
    \item Estados T (tenso) e R (relaxado)
    \item Transição concertada
    \item Todas subunidades no mesmo estado
    \item Ligante estabiliza estado R
\end{itemize}
\end{block}

\begin{center}
\begin{tikzpicture}[scale=0.5]
    % T state
    \foreach \i in {0,1,2,3}
        \draw[thick, fill=red!30] ({cos(90*\i)},{sin(90*\i)}) circle (0.3);
    \node[below] at (0,-1.5) {\tiny T (baixa afinidade)};

    % Arrow
    \draw[->, thick] (2,0) -- (3,0);
    \node[above, font=\tiny] at (2.5,0.2) {$+S$};

    % R state
    \begin{scope}[xshift=5.5cm]
        \foreach \i in {0,1,2,3}
            \draw[thick, fill=verdeacido!50] ({cos(90*\i)},{sin(90*\i)}) circle (0.3);
        \node[below] at (0,-1.5) {\tiny R (alta afinidade)};
    \end{scope}
\end{tikzpicture}
\end{center}

\column{0.5\textwidth}
\begin{block}{Modelo KNF}
\textbf{Koshland-Némethy-Filmer}
\begin{itemize}
    \item Transição sequencial
    \item Subunidades mudam individualmente
    \item Ligação induz mudança conformacional
    \item Afeta subunidades vizinhas
\end{itemize}
\end{block}

\begin{center}
\begin{tikzpicture}[scale=0.5]
    % Initial state
    \foreach \i in {0,1,2,3}
        \draw[thick, fill=red!30] ({cos(90*\i)},{sin(90*\i)}) circle (0.3);

    % Arrow
    \draw[->, thick] (2,0) -- (3,0);
    \node[above, font=\tiny] at (2.5,0.2) {$+S$};

    % Intermediate state
    \begin{scope}[xshift=5.5cm]
        \draw[thick, fill=verdeacido!50] (0,1) circle (0.3);
        \draw[thick, fill=verdeacido!50] (1,0) circle (0.3);
        \draw[thick, fill=red!30] (0,-1) circle (0.3);
        \draw[thick, fill=red!30] (-1,0) circle (0.3);
    \end{scope}
\end{tikzpicture}
\end{center}
\end{columns}
\end{frame}

\begin{frame}{Inibição Enzimática: Tipos}
\begin{block}{Inibição Competitiva}
Inibidor compete com substrato pelo sítio ativo
\begin{equation*}
v = \frac{V_{max}[S]}{K_M\left(1 + \frac{[I]}{K_I}\right) + [S]}
\end{equation*}

$V_{max}$ inalterado, $K_M^{app}$ aumenta
\end{block}

\begin{block}{Inibição Não-Competitiva}
Inibidor liga em sítio diferente (não afeta ligação de S)
\begin{equation*}
v = \frac{V_{max}[S]}{\left(1 + \frac{[I]}{K_I}\right)(K_M + [S])}
\end{equation*}

$K_M$ inalterado, $V_{max}^{app}$ diminui
\end{block}
\end{frame}

\begin{frame}{Inibição Enzimática: Tipos Adicionais}
\begin{block}{Inibição Incompetitiva}
Inibidor liga apenas ao complexo ES
\begin{equation*}
v = \frac{V_{max}[S]}{K_M + [S]\left(1 + \frac{[I]}{K_I}\right)}
\end{equation*}

$K_M$ e $V_{max}$ ambos diminuem proporcionalmente
\end{block}

\begin{block}{Inibição Mista}
Inibidor pode ligar tanto a E quanto a ES, com afinidades diferentes
\begin{equation*}
v = \frac{V_{max}[S]}{K_M\left(1 + \frac{[I]}{K_I}\right) + [S]\left(1 + \frac{[I]}{K_I'}\right)}
\end{equation*}

Afeta tanto $K_M$ quanto $V_{max}$
\end{block}
\end{frame}

\begin{frame}{Lineweaver-Burk: Identificando Tipo de Inibição}
\begin{center}
\begin{tikzpicture}[scale=0.65]
    % Competitive
    \begin{scope}
        \draw[->, thick] (-1.5,0) -- (3,0) node[right, font=\tiny] {$1/[S]$};
        \draw[->, thick] (0,-0.3) -- (0,3) node[above, font=\tiny] {$1/v$};
        \draw[azulescuro, thick, domain=-1:2.5] plot (\x, {0.5*\x + 1});
        \draw[red, thick, domain=-0.5:2.5] plot (\x, {0.8*\x + 1});
        \node[below, font=\tiny] at (1,-0.8) {Competitiva};
        \node[right, font=\tiny, azulescuro] at (2.5,2.2) {Sem I};
        \node[right, font=\tiny, red] at (2.5,3) {Com I};
    \end{scope}

    % Non-competitive
    \begin{scope}[xshift=5.5cm]
        \draw[->, thick] (-1.5,0) -- (3,0) node[right, font=\tiny] {$1/[S]$};
        \draw[->, thick] (0,-0.3) -- (0,3) node[above, font=\tiny] {$1/v$};
        \draw[azulescuro, thick, domain=-1.2:2.5] plot (\x, {0.5*\x + 1});
        \draw[red, thick, domain=-1.2:2.5] plot (\x, {0.5*\x + 1.5});
        \node[below, font=\tiny] at (1,-0.8) {Não-Competitiva};
    \end{scope}

    % Uncompetitive
    \begin{scope}[yshift=-4.5cm]
        \draw[->, thick] (-1.5,0) -- (3,0) node[right, font=\tiny] {$1/[S]$};
        \draw[->, thick] (0,-0.3) -- (0,3) node[above, font=\tiny] {$1/v$};
        \draw[azulescuro, thick, domain=-1.2:2.5] plot (\x, {0.5*\x + 1});
        \draw[red, thick, domain=-0.7:2.5] plot (\x, {0.5*\x + 1.5});
        \node[below, font=\tiny] at (1,-0.8) {Incompetitiva};
    \end{scope}

    % Mixed
    \begin{scope}[xshift=5.5cm, yshift=-4.5cm]
        \draw[->, thick] (-1.5,0) -- (3,0) node[right, font=\tiny] {$1/[S]$};
        \draw[->, thick] (0,-0.3) -- (0,3) node[above, font=\tiny] {$1/v$};
        \draw[azulescuro, thick, domain=-1.2:2.5] plot (\x, {0.5*\x + 1});
        \draw[red, thick, domain=-0.8:2.5] plot (\x, {0.7*\x + 1.4});
        \node[below, font=\tiny] at (1,-0.8) {Mista};
    \end{scope}
\end{tikzpicture}
\end{center}
\end{frame}

\begin{frame}{Inibição Reversível vs Irreversível}
\begin{columns}
\column{0.5\textwidth}
\begin{block}{Reversível}
Inibidor liga não-covalentemente
\begin{itemize}
    \item Equilíbrio dinâmico
    \item Pode ser removido (diálise)
    \item Constante de dissociação $K_I$
    \item Exemplos: metabólitos regulatórios
\end{itemize}
\end{block}

\exemplo{Estatinas}{
Inibidores competitivos reversíveis da HMG-CoA redutase (síntese de colesterol)
}

\column{0.5\textwidth}
\begin{block}{Irreversível}
Inibidor forma ligação covalente
\begin{itemize}
    \item Inativação permanente
    \item Requer síntese de nova enzima
    \item Taxa de inativação $k_{inact}$
    \item Exemplos: toxinas, fármacos
\end{itemize}
\end{block}

\exemplo{Aspirina}{
Acetila irreversivelmente Ser530 da COX (inibe inflamação)
}
\end{columns}

\vspace{0.3cm}

\importante{
$IC_{50}$: concentração de inibidor que reduz atividade em 50\%\\
$K_I$: constante de dissociação do complexo EI (afinidade)
}
\end{frame}

\begin{frame}[fragile]{Ajuste de Dados de Cinética Enzimática}
\begin{lstlisting}[language=Python, caption=Regressão não-linear para Michaelis-Menten, basicstyle=\ttfamily\footnotesize]
import numpy as np
from scipy.optimize import curve_fit
import matplotlib.pyplot as plt

# Funcao de Michaelis-Menten
def michaelis_menten(S, Vmax, Km):
    return Vmax * S / (Km + S)

# Dados experimentais
S = np.array([0.5, 1, 2, 5, 10, 20, 50])  # mM
v = np.array([15, 25, 38, 60, 75, 85, 95])  # umol/min

# Ajuste nao-linear
params, covariance = curve_fit(michaelis_menten, S, v,
                               p0=[100, 5])  # estimativas iniciais

Vmax_fit, Km_fit = params
print(f"Vmax = {Vmax_fit:.2f} umol/min")
print(f"Km = {Km_fit:.2f} mM")
print(f"kcat/Km = {Vmax_fit/Km_fit:.2f} min^-1 mM^-1")

# Plotar
S_fit = np.linspace(0, 50, 200)
v_fit = michaelis_menten(S_fit, Vmax_fit, Km_fit)
plt.plot(S, v, 'o', label='Dados')
plt.plot(S_fit, v_fit, '-', label='Ajuste')
plt.xlabel('[S] (mM)')
plt.ylabel('v (umol/min)')
plt.legend()
\end{lstlisting}
\end{frame}

\begin{frame}[fragile]{Cálculo de IC50 e Ki}
\begin{lstlisting}[language=Python, caption=Determinação de IC50, basicstyle=\ttfamily\footnotesize]
from scipy.optimize import curve_fit

# Funcao dose-resposta (Hill)
def dose_response(I, IC50, n):
    """Inibicao em funcao de [I]"""
    return 100 / (1 + (I/IC50)**n)

# Dados de inibicao
I = np.array([0.01, 0.1, 1, 10, 100, 1000])  # uM
activity = np.array([98, 90, 70, 30, 8, 2])  # % atividade

# Ajuste
params, _ = curve_fit(dose_response, I, activity, p0=[10, 1])
IC50, n = params

print(f"IC50 = {IC50:.2f} uM")
print(f"Coeficiente de Hill = {n:.2f}")

# Para inibidor competitivo, calcular Ki
S = 5  # mM (concentracao de substrato usada)
Km = 2  # mM
Ki = IC50 / (1 + S/Km)
print(f"Ki = {Ki:.2f} uM")

# Plotar curva dose-resposta
I_fit = np.logspace(-2, 3, 200)
activity_fit = dose_response(I_fit, IC50, n)
plt.semilogx(I, activity, 'o', label='Dados')
plt.semilogx(I_fit, activity_fit, '-', label='Ajuste')
plt.axhline(y=50, color='r', linestyle='--', label='IC50')
plt.axvline(x=IC50, color='r', linestyle='--')
plt.xlabel('[Inibidor] (uM)')
plt.ylabel('Atividade (%)')
plt.legend()
\end{lstlisting}
\end{frame}

% ============================================================================
% SEÇÃO 4: CASCATAS DE SINALIZAÇÃO
% ============================================================================
\section{Cascatas de Sinalização Celular}

\begin{frame}{7.3 Sinalização Celular: Visão Geral}
\definicao{Transdução de Sinal}{
Processo pelo qual uma célula converte um sinal extracelular em uma resposta intracelular específica.
}

\vspace{0.3cm}

\begin{center}
\begin{tikzpicture}[node distance=2cm, auto,
    box/.style={rectangle, draw, text width=2.5cm, text centered, rounded corners, minimum height=0.8cm, font=\small}]

    \node[box, fill=roxodna!30] (lig) {Ligante\\(hormônio, GF)};
    \node[box, fill=laranjacel!30, below of=lig] (rec) {Receptor\\(membrana)};
    \node[box, fill=verdeacido!30, below of=rec] (trans) {Transdutores\\(proteínas G, quinases)};
    \node[box, fill=azulclaro!30, below of=trans] (eff) {Efetores\\(enzimas, TFs)};
    \node[box, fill=yellow!30, below of=eff] (resp) {Resposta\\(gene, metabolismo)};

    \draw[->, thick] (lig) -- (rec);
    \draw[->, thick] (rec) -- (trans);
    \draw[->, thick] (trans) -- (eff);
    \draw[->, thick] (eff) -- (resp);
\end{tikzpicture}
\end{center}
\end{frame}

\begin{frame}{Tipos de Receptores de Membrana}
\begin{columns}
\column{0.5\textwidth}
\begin{block}{GPCRs}
\textbf{Receptores Acoplados a Proteína G}
\begin{itemize}
    \item 7 hélices transmembrana
    \item Ligação ativa proteína G
    \item Maior família de receptores
    \item Alvos de 30-40\% dos fármacos
\end{itemize}
Exemplos: receptores adrenérgicos, dopaminérgicos, olfativos
\end{block}

\begin{block}{RTKs}
\textbf{Receptores Tirosina Quinase}
\begin{itemize}
    \item 1 hélice transmembrana
    \item Atividade quinase intrínseca
    \item Dimerização por ligante
    \item Fosforilação de tirosinas
\end{itemize}
Exemplos: EGFR, VEGFR, insulina
\end{block}

\column{0.5\textwidth}
\begin{block}{Canais Iônicos}
\textbf{Receptores Ionotrópicos}
\begin{itemize}
    \item Ligante abre canal
    \item Fluxo de íons (Na$^+$, K$^+$, Ca$^{2+}$, Cl$^-$)
    \item Resposta rápida (ms)
\end{itemize}
Exemplos: receptor nicotínico, GABA$_A$
\end{block}

\begin{block}{Receptores Nucleares}
\textbf{Fatores de Transcrição}
\begin{itemize}
    \item Intracelulares
    \item Ligante lipofílico
    \item Regulação direta de genes
\end{itemize}
Exemplos: receptores de esteroides, tireoide
\end{block}
\end{columns}
\end{frame}

\begin{frame}{Via GPCR: Proteínas G Heterotriméricas}
\begin{center}
\begin{tikzpicture}[scale=0.7]
    % Membrane
    \draw[thick, double] (-4,0) -- (4,0);
    \node[left, font=\tiny] at (-4,0.2) {Membrana};

    % GPCR
    \draw[thick, fill=roxodna!30] (-2,-0.5) rectangle (-1,0.5);
    \node[font=\tiny] at (-1.5,0) {GPCR};

    % Ligand
    \node[draw, circle, fill=laranjacel!50, minimum size=0.4cm, font=\tiny] at (-1.5,1.2) {L};
    \draw[->, thick] (-1.5,0.9) -- (-1.5,0.6);

    % G protein (inactive)
    \node[draw, circle, fill=azulclaro!30, font=\tiny] at (0,-1) {$\alpha$};
    \node[font=\tiny] at (0,-1.3) {GDP};
    \node[draw, ellipse, fill=verdeacido!30, font=\tiny] at (1,-1) {$\beta\gamma$};

    % Arrow
    \draw[->, ultra thick, azulescuro] (2,-1) -- (3.5,-1);
    \node[above, font=\tiny] at (2.75,-0.8) {Ativação};

    % G protein (active)
    \begin{scope}[xshift=5.5cm]
        \node[draw, circle, fill=red!50, font=\tiny] at (0,-1) {$\alpha$};
        \node[font=\tiny] at (0,-1.3) {GTP};

        \node[draw, ellipse, fill=verdeacido!30, font=\tiny] at (1.5,-1) {$\beta\gamma$};

        \draw[->, thick] (0,-1.5) -- (0,-2.5);
        \node[below, font=\tiny, text width=2cm, align=center] at (0,-2.8) {Ativa\\efetores};
    \end{scope}
\end{tikzpicture}
\end{center}

\vspace{0.3cm}

\begin{columns}
\column{0.5\textwidth}
\textbf{Subunidades G$\alpha$:}
\begin{itemize}
    \item G$\alpha_s$: estimula adenilil ciclase
    \item G$\alpha_i$: inibe adenilil ciclase
    \item G$\alpha_q$: ativa fosfolipase C
\end{itemize}

\column{0.5\textwidth}
\textbf{Subunidade G$\beta\gamma$:}
\begin{itemize}
    \item Ativa canais iônicos
    \item Regula adenilil ciclase
    \item Ativa PI3K
\end{itemize}
\end{columns}
\end{frame}

\begin{frame}{Segundos Mensageiros}
\begin{columns}
\column{0.5\textwidth}
\begin{block}{cAMP}
\textbf{AMP cíclico}
\begin{itemize}
    \item Produzido por adenilil ciclase
    \item Ativa PKA (proteína quinase A)
    \item Degradado por fosfodiesterases
    \item Via G$\alpha_s$
\end{itemize}
\end{block}

\begin{block}{Ca$^{2+}$}
\textbf{Cálcio intracelular}
\begin{itemize}
    \item Liberado do RE
    \item Liga calmodulina
    \item Ativa CaMK, PKC
    \item Concentração: 100 nM $\rightarrow$ 1 $\mu$M
\end{itemize}
\end{block}

\column{0.5\textwidth}
\begin{block}{IP$_3$ e DAG}
\textbf{Inositol-trifosfato e Diacilglicerol}
\begin{itemize}
    \item Produzidos por fosfolipase C
    \item IP$_3$ libera Ca$^{2+}$ do RE
    \item DAG ativa PKC
    \item Via G$\alpha_q$
\end{itemize}
\end{block}

\begin{block}{cGMP}
\textbf{GMP cíclico}
\begin{itemize}
    \item Produzido por guanilil ciclase
    \item Ativa PKG
    \item Via óxido nítrico (NO)
    \item Vasodilatação
\end{itemize}
\end{block}
\end{columns}
\end{frame}

\begin{frame}{Cascata MAPK: Via Clássica}
\begin{center}
\begin{tikzpicture}[node distance=1.5cm, auto,
    box/.style={rectangle, draw, text width=2.2cm, text centered, rounded corners, minimum height=0.7cm, font=\small}]

    \node[box, fill=roxodna!30] (gf) {Growth Factor};
    \node[box, fill=laranjacel!30, below of=gf] (rtk) {RTK};
    \node[box, fill=verdeacido!30, below of=rtk] (ras) {Ras-GTP};
    \node[box, fill=azulclaro!30, below of=ras] (raf) {Raf (MAPKKK)};
    \node[box, fill=yellow!30, below of=raf] (mek) {MEK (MAPKK)};
    \node[box, fill=red!20, below of=mek] (erk) {ERK (MAPK)};
    \node[box, fill=verdeacido!50, below of=erk] (tf) {Fatores de\\Transcrição};

    \draw[->, thick] (gf) -- (rtk);
    \draw[->, thick] (rtk) -- node[right, font=\tiny] {Grb2/SOS} (ras);
    \draw[->, thick] (ras) -- (raf);
    \draw[->, thick] (raf) -- node[right, font=\tiny] {P} (mek);
    \draw[->, thick] (mek) -- node[right, font=\tiny] {P P} (erk);
    \draw[->, thick] (erk) -- node[right, font=\tiny] {P} (tf);

    % Nucleus
    \draw[thick, dashed] (-2.5,-7.5) rectangle (2.5,-9);
    \node[font=\tiny] at (0,-7.2) {Núcleo};
    \node[box, fill=roxodna!30, minimum width=1.5cm] at (0,-8.25) {DNA};
    \draw[->, thick] (tf) -- (0,-7.5);
\end{tikzpicture}
\end{center}
\end{frame}

\begin{frame}{MAPK: Módulo de Amplificação}
\importante{
Cascata de fosforilação amplifica sinal:
\begin{itemize}
    \item 1 RTK ativa múltiplas moléculas de Ras
    \item 1 Ras ativa múltiplas moléculas de Raf
    \item 1 Raf fosforila múltiplas MEK
    \item 1 MEK fosforila múltiplas ERK
\end{itemize}
}

\vspace{0.3cm}

\begin{block}{Amplificação em Cascata}
Se cada quinase ativa 10 moléculas da próxima:
\begin{equation*}
\text{Amplificação} = 10^3 = 1000\times
\end{equation*}
\end{block}

\vspace{0.3cm}

\begin{exampleblock}{Outras Vias MAPK}
\begin{itemize}
    \item \textbf{JNK/SAPK:} estresse, apoptose
    \item \textbf{p38 MAPK:} inflamação, diferenciação
\end{itemize}
\end{exampleblock}
\end{frame}

\begin{frame}{Dinâmica de Fosforilação e Desfosforilação}
\begin{block}{Ciclo de Fosforilação}
\begin{equation*}
\text{Proteína} + \text{ATP} \xrightarrow{\text{Quinase}} \text{Proteína-P} + \text{ADP}
\end{equation*}
\begin{equation*}
\text{Proteína-P} + \text{H}_2\text{O} \xrightarrow{\text{Fosfatase}} \text{Proteína} + \text{P}_i
\end{equation*}
\end{block}

\vspace{0.3cm}

\begin{columns}
\column{0.5\textwidth}
\textbf{Quinases:}
\begin{itemize}
    \item Fosforilam Ser, Thr, Tyr
    \item $\sim$518 quinases humanas
    \item Domínio catalítico conservado
    \item Regulação por fosforilação
\end{itemize}

\column{0.5\textwidth}
\textbf{Fosfatases:}
\begin{itemize}
    \item Removem fosfatos
    \item Menos específicas
    \item Regulação por subunidades
    \item Antagonizam quinases
\end{itemize}
\end{columns}

\vspace{0.3cm}

\destaque{
Balanço quinase/fosfatase determina estado de fosforilação
}
\end{frame}

\begin{frame}{Ultrassensibilidade de Goldbeter-Koshland}
\begin{columns}
\column{0.5\textwidth}
\begin{block}{Modelo Matemático}
Quando quinase e fosfatase operam próximas à saturação:

\begin{equation*}
\frac{[\text{Proteína-P}]}{[\text{Proteína}]_{total}} = \frac{v_1 - v_2}{v_1 + v_2}
\end{equation*}

Resposta switch-like mesmo sem cooperatividade!
\end{block}

\textbf{Condições:}
\begin{itemize}
    \item Alta [Proteína] vs. $K_M$
    \item Ciclo de fosforilação/desfosforilação
    \item Saturação de ambas enzimas
\end{itemize}

\column{0.5\textwidth}
\begin{center}
\begin{tikzpicture}[scale=0.7]
    % Axes
    \draw[->, thick] (0,0) -- (5,0) node[right, font=\small] {Sinal};
    \draw[->, thick] (0,0) -- (0,4) node[above, font=\small] {Resposta};

    % Michaelian (hyperbolic)
    \draw[thick, azulescuro, domain=0:4.8, samples=100] plot (\x, {3*\x/(1+\x)});
    \node[right, font=\tiny, text width=2cm] at (3,2) {Michaeliana\\(gradual)};

    % Ultrasensitive (sigmoidal)
    \draw[thick, red, domain=0:4.8, samples=100] plot (\x, {3*\x^4/(1+\x^4)});
    \node[right, font=\tiny, text width=2cm] at (1.5,1) {Ultrassensível\\(switch-like)};

    % Threshold
    \draw[dashed, verdeacido] (1.5,0) -- (1.5,3.5);
    \node[above, font=\tiny] at (1.5,3.7) {Limiar};
\end{tikzpicture}
\end{center}

\importante{
Permite decisões digitais (tudo-ou-nada) em sistemas analógicos
}
\end{columns}
\end{frame}

\begin{frame}[fragile]{Modelagem de Cascatas de Sinalização}
\begin{lstlisting}[language=Python, caption=Cascata MAPK simplificada, basicstyle=\ttfamily\footnotesize]
import numpy as np
from scipy.integrate import odeint
import matplotlib.pyplot as plt

def mapk_cascade(y, t, k1, k2, k3):
    """
    Modelo simplificado de cascata MAPK
    y = [MAPKKK*, MAPKK*, MAPK*] (* = fosforilado/ativo)
    """
    MAPKKK_active, MAPKK_active, MAPK_active = y

    # Assumindo input constante e desfosforilacao linear
    dMAPKKK = k1 - 0.1*MAPKKK_active
    dMAPKK = k2*MAPKKK_active - 0.1*MAPKK_active
    dMAPK = k3*MAPKK_active - 0.1*MAPK_active

    return [dMAPKKK, dMAPKK, dMAPK]

# Parametros
k1, k2, k3 = 1.0, 2.0, 3.0  # taxas de ativacao
y0 = [0, 0, 0]  # condicoes iniciais
t = np.linspace(0, 50, 500)

# Resolver EDO
solution = odeint(mapk_cascade, y0, t, args=(k1, k2, k3))

# Plotar
plt.plot(t, solution[:, 0], label='MAPKKK*')
plt.plot(t, solution[:, 1], label='MAPKK*')
plt.plot(t, solution[:, 2], label='MAPK*', linewidth=2)
plt.xlabel('Tempo (min)')
plt.ylabel('Concentracao (uM)')
plt.title('Cascata MAPK: Amplificacao de Sinal')
plt.legend()
plt.grid(True)
\end{lstlisting}
\end{frame}

% ============================================================================
% SEÇÃO 5: MODIFICAÇÕES PÓS-TRADUCIONAIS
% ============================================================================
\section{Modificações Pós-Traducionais}

\begin{frame}{7.4 Modificações Pós-Traducionais (PTMs)}
\definicao{PTM}{
Modificações químicas covalentes em proteínas após tradução, expandindo diversidade e regulação do proteoma.
}

\vspace{0.3cm}

\begin{columns}
\column{0.5\textwidth}
\textbf{Tipos Principais:}
\begin{itemize}
    \item Fosforilação (Ser, Thr, Tyr)
    \item Acetilação (Lys)
    \item Metilação (Lys, Arg)
    \item Ubiquitinação (Lys)
    \item Glicosilação (Asn, Ser, Thr)
    \item SUMOilação (Lys)
    \item Nitrosilação (Cys)
    \item Palmitoilação (Cys)
\end{itemize}

\column{0.5\textwidth}
\importante{
PTMs regulam:
\begin{itemize}
    \item Atividade enzimática
    \item Localização subcelular
    \item Interações proteína-proteína
    \item Estabilidade
    \item Conformação
\end{itemize}
}

\vspace{0.3cm}

\begin{alertblock}{Estatística}
$>$400 tipos de PTMs conhecidas!\\
Proteína média: 5-10 PTMs diferentes
\end{alertblock}
\end{columns}
\end{frame}

\begin{frame}{Fosforilação: PTM mais Estudada}
\begin{block}{Características}
\begin{itemize}
    \item Mais abundante e dinâmica
    \item Reversível (quinases/fosfatases)
    \item Adição de grupo fosfato (PO$_4^{3-}$)
    \item Carga negativa: muda interações eletrostáticas
    \item $\sim$30\% das proteínas fosforiladas a qualquer momento
\end{itemize}
\end{block}

\vspace{0.3cm}

\begin{columns}
\column{0.5\textwidth}
\textbf{Resíduos Alvo:}
\begin{itemize}
    \item \textbf{Ser:} 90\% (Ser-OH)
    \item \textbf{Thr:} 10\% (Thr-OH)
    \item \textbf{Tyr:} $<$1\% (Tyr-OH)
    \item His, Asp (procariotos)
\end{itemize}

\column{0.5\textwidth}
\textbf{Consequências:}
\begin{itemize}
    \item Ativação/inativação enzimática
    \item Criação de sítios de ligação (SH2, 14-3-3)
    \item Mudanças conformacionais
    \item Marcação para degradação
\end{itemize}
\end{columns}
\end{frame}

\begin{frame}{Acetilação e Metilação: Regulação Epigenética}
\begin{columns}
\column{0.5\textwidth}
\begin{block}{Acetilação}
\textbf{Adição de grupo acetil (COCH$_3$)}
\begin{itemize}
    \item Principalmente em Lys
    \item Neutraliza carga positiva
    \item HATs: histona acetiltransferases
    \item HDACs: histona deacetilases
\end{itemize}

\textbf{Efeitos em Histonas:}
\begin{itemize}
    \item Reduz afinidade DNA-histona
    \item Cromatina aberta (eucromatina)
    \item Ativa transcrição
\end{itemize}
\end{block}

\column{0.5\textwidth}
\begin{block}{Metilação}
\textbf{Adição de grupos metil (CH$_3$)}
\begin{itemize}
    \item Lys (mono, di, tri-metil)
    \item Arg (mono, di-metil)
    \item HMTs: histona metiltransferases
    \item HDMs: histona demetilases
\end{itemize}

\textbf{Efeitos Contextuais:}
\begin{itemize}
    \item H3K4me3: ativação
    \item H3K9me3: repressão
    \item H3K27me3: silenciamento
\end{itemize}
\end{block}
\end{columns}

\vspace{0.3cm}

\destaque{
Código de histonas: combinação de PTMs determina estado transcricional
}
\end{frame}

\begin{frame}{Ubiquitinação e Degradação Proteica}
\begin{block}{Sistema Ubiquitina-Proteassoma}
\textbf{Ubiquitina:} proteína pequena (76 aa) que marca alvos para degradação
\end{block}

\vspace{0.3cm}

\begin{center}
\begin{tikzpicture}[node distance=1.3cm, auto,
    box/.style={rectangle, draw, text width=2cm, text centered, rounded corners, minimum height=0.6cm, font=\small}]

    \node[box, fill=azulclaro!30] (e1) {E1\\(ativação)};
    \node[box, fill=verdeacido!30, right of=e1] (e2) {E2\\(conjugação)};
    \node[box, fill=laranjacel!30, right of=e2] (e3) {E3\\(ligação)};
    \node[box, fill=roxodna!30, below of=e3] (target) {Proteína\\Alvo};
    \node[box, fill=red!30, below of=e1] (protea) {Proteassoma\\26S};

    \draw[->, thick] (e1) -- node[above, font=\tiny] {Ub} (e2);
    \draw[->, thick] (e2) -- node[above, font=\tiny] {Ub} (e3);
    \draw[->, thick] (e3) -- node[right, font=\tiny] {Ub$_n$} (target);
    \draw[->, thick] (target) -- node[above, font=\tiny] {Ub$_4$-Proteína} (protea);

    \node[below, font=\tiny, text width=3cm, align=center] at (0,-2.5) {Degradação\\a aminoácidos};
\end{tikzpicture}
\end{center}

\vspace{0.3cm}

\begin{columns}
\column{0.5\textwidth}
\textbf{Tipos de Cadeia:}
\begin{itemize}
    \item K48: degradação (26S)
    \item K63: sinalização (não-degradativa)
    \item Mono-Ub: endocitose, reparo DNA
\end{itemize}

\column{0.5\textwidth}
\importante{
E3 ligases conferem especificidade:\\
$>$600 E3s humanas!
}
\end{columns}
\end{frame}

\begin{frame}{Glicosilação: Modificação Mais Complexa}
\begin{columns}
\column{0.5\textwidth}
\begin{block}{N-Glicosilação}
\textbf{Em Asn (Asn-X-Ser/Thr)}
\begin{itemize}
    \item Ocorre no RE
    \item Oligossacarídeo pré-formado
    \item Transferido en bloc
    \item Processamento no Golgi
\end{itemize}
\end{block}

\begin{block}{O-Glicosilação}
\textbf{Em Ser/Thr}
\begin{itemize}
    \item Construção sequencial
    \item Golgi e citoplasma
    \item Mais diversa
    \item O-GlcNAc (citoplasma)
\end{itemize}
\end{block}

\column{0.5\textwidth}
\textbf{Funções:}
\begin{itemize}
    \item Dobramento proteico
    \item Estabilidade
    \item Reconhecimento celular
    \item Adesão celular
    \item Imunidade (anticorpos)
\end{itemize}

\vspace{0.3cm}

\begin{exampleblock}{Importância Clínica}
\begin{itemize}
    \item Hemoglobina glicosilada (HbA1c): diabetes
    \item Glicosilação aberrante: câncer
    \item Grupo sanguíneo: glicosilação de RBCs
\end{itemize}
\end{exampleblock}
\end{columns}
\end{frame}

\begin{frame}{Crosstalk entre PTMs}
\importante{
PTMs não agem isoladamente: interagem e se regulam mutuamente!
}

\vspace{0.3cm}

\begin{block}{Exemplos de Crosstalk}
\begin{itemize}
    \item \textbf{Fosforilação inibe acetilação:} mesmo resíduo (Ser/Thr vs. Lys adjacente)
    \item \textbf{Ubiquitinação requer fosforilação:} fosfo-degron
    \item \textbf{O-GlcNAcilação compete com fosforilação:} mesmas Ser/Thr
    \item \textbf{Metilação regula fosforilação:} em histonas
    \item \textbf{SUMOilação antagoniza ubiquitinação:} mesmas Lys
\end{itemize}
\end{block}

\vspace{0.3cm}

\exemplo{p53}{
Acetilação (ativa) vs. Ubiquitinação (degrada) em Lys$_{382}$ regula estabilidade do supressor tumoral p53
}
\end{frame}

\begin{frame}{Bancos de Dados de PTMs}
\begin{columns}
\column{0.5\textwidth}
\begin{block}{PhosphoSitePlus}
\url{https://www.phosphosite.org}
\begin{itemize}
    \item PTMs experimentalmente verificadas
    \item $>$500,000 sítios
    \item Fosforilação, acetilação, ubiquitinação, etc.
    \item Regulação up/down
    \item Informação funcional
\end{itemize}
\end{block}

\begin{block}{UniProt}
\url{https://www.uniprot.org}
\begin{itemize}
    \item PTMs anotadas
    \item Integração com literatura
    \item Variantes e mutações
\end{itemize}
\end{block}

\column{0.5\textwidth}
\begin{block}{Outros Recursos}
\begin{itemize}
    \item \textbf{dbPTM:} integração de PTMs
    \item \textbf{PTMcode:} crosstalk
    \item \textbf{PHOSIDA:} fosforilação
    \item \textbf{PLMD:} metilação de lisinas
    \item \textbf{O-GlcNAcAtlas:} O-GlcNAc
\end{itemize}
\end{block}

\vspace{0.3cm}

\begin{exampleblock}{Detecção}
\textbf{Técnica Principal:}\\
Espectrometria de massa (LC-MS/MS)
\begin{itemize}
    \item Identifica modificações
    \item Mapeia sítios exatos
    \item Quantificação relativa/absoluta
\end{itemize}
\end{exampleblock}
\end{columns}
\end{frame}

\begin{frame}{PTMs no Controle do Ciclo Celular}
\exemplo{Fosforilação de Ciclinas/CDKs}{
Progressão pelo ciclo celular é regulada por fosforilação:
}

\begin{center}
\begin{tikzpicture}[scale=0.7]
    % Cell cycle
    \draw[thick] (0,0) circle (2);

    % Phases
    \node at (-1.4,1.4) {G1};
    \node at (1.4,1.4) {S};
    \node at (1.4,-1.4) {G2};
    \node at (-1.4,-1.4) {M};

    % CDK/Cyclin complexes
    \node[draw, fill=verdeacido!30, rounded corners, font=\tiny] at (-1,2.5) {Cyclin D/CDK4,6};
    \node[draw, fill=azulclaro!30, rounded corners, font=\tiny] at (2.5,1) {Cyclin E/CDK2};
    \node[draw, fill=laranjacel!30, rounded corners, font=\tiny] at (2.5,-1) {Cyclin A/CDK2};
    \node[draw, fill=roxodna!30, rounded corners, font=\tiny] at (-1,-2.5) {Cyclin B/CDK1};

    % Arrows
    \draw[->, thick] (-1,2.3) -- (-1,1.7);
    \draw[->, thick] (2.3,1) -- (1.7,1);
    \draw[->, thick] (2.3,-1) -- (1.7,-1);
    \draw[->, thick] (-1,-2.3) -- (-1,-1.7);

    % Checkpoints
    \foreach \angle in {45, 135, 225, 315} {
        \fill[red] ({2*cos(\angle)}, {2*sin(\angle)}) circle (2pt);
    }
\end{tikzpicture}
\end{center}

\vspace{0.3cm}

\begin{itemize}
    \item \textbf{Fosforilação de Rb:} libera E2F $\rightarrow$ transição G1/S
    \item \textbf{Fosforilação de histonas:} condensação cromossomal
    \item \textbf{Desfosforilação de CDK1:} entrada em mitose
\end{itemize}
\end{frame}

% ============================================================================
% SEÇÃO 6: PROTEÔMICA
% ============================================================================
\section{Proteômica}

\begin{frame}{7.5 Proteômica: Estudo do Proteoma}
\definicao{Proteoma}{
Conjunto completo de proteínas expressas por um genoma, célula ou tecido em determinado momento e condição.
}

\vspace{0.3cm}

\begin{columns}
\column{0.5\textwidth}
\textbf{Por que Proteômica?}
\begin{itemize}
    \item mRNA $\neq$ Proteína
    \item PTMs não previstas por DNA
    \item Abundância dinâmica
    \item Interações proteína-proteína
    \item Localização subcelular
\end{itemize}

\column{0.5\textwidth}
\textbf{Desafios:}
\begin{itemize}
    \item Faixa dinâmica enorme ($10^6$)
    \item Diversidade química
    \item Modificações pós-traducionais
    \item Complexos transitórios
    \item Sem amplificação (vs. PCR)
\end{itemize}
\end{columns}

\vspace{0.3cm}

\importante{
Proteoma humano: $\sim$20,000 genes $\rightarrow$ $>$1,000,000 proteoformas (com PTMs e isoformas)
}
\end{frame}

\begin{frame}{Abordagens Proteômicas}
\begin{columns}
\column{0.5\textwidth}
\begin{block}{Bottom-Up}
\textbf{Shotgun Proteomics}
\begin{enumerate}
    \item Digestão proteolítica (tripsina)
    \item Peptídeos analisados por MS
    \item Identificação por busca em banco
    \item Reconstrução de proteínas
\end{enumerate}

\textbf{Vantagens:}
\begin{itemize}
    \item Mais sensível
    \item Cobertura ampla
    \item PTMs mapeadas
\end{itemize}
\end{block}

\column{0.5\textwidth}
\begin{block}{Top-Down}
\textbf{Proteínas Intactas}
\begin{enumerate}
    \item Proteínas inteiras no MS
    \item Fragmentação no espectrômetro
    \item Análise direta
\end{enumerate}

\textbf{Vantagens:}
\begin{itemize}
    \item Isoformas preservadas
    \item Combinações de PTMs
    \item Sem ambiguidade de montagem
\end{itemize}

\textbf{Desvantagens:}
\begin{itemize}
    \item Menos sensível
    \item Proteínas grandes difíceis
\end{itemize}
\end{block}
\end{columns}
\end{frame}

\begin{frame}{Espectrometria de Massa: LC-MS/MS}
\begin{center}
\begin{tikzpicture}[node distance=2cm, auto,
    box/.style={rectangle, draw, text width=2cm, text centered, rounded corners, minimum height=0.7cm, font=\small}]

    \node[box, fill=roxodna!30] (sample) {Amostra\\Proteica};
    \node[box, fill=laranjacel!30, below of=sample] (digest) {Digestão\\(Tripsina)};
    \node[box, fill=verdeacido!30, below of=digest] (lc) {LC\\(Separação)};
    \node[box, fill=azulclaro!30, below of=lc] (ms1) {MS1\\(m/z)};
    \node[box, fill=yellow!30, below of=ms1] (ms2) {MS2\\(Fragmentação)};
    \node[box, fill=red!20, below of=ms2] (id) {Identificação\\(Banco dados)};

    \draw[->, thick] (sample) -- (digest);
    \draw[->, thick] (digest) -- node[right, font=\tiny] {Peptídeos} (lc);
    \draw[->, thick] (lc) -- (ms1);
    \draw[->, thick] (ms1) -- node[right, font=\tiny] {Precursor} (ms2);
    \draw[->, thick] (ms2) -- node[right, font=\tiny] {Espectro} (id);
\end{tikzpicture}
\end{center}

\vspace{0.3cm}

\begin{columns}
\column{0.5\textwidth}
\textbf{MS1:}
\begin{itemize}
    \item Massa/carga (m/z)
    \item Intensidade do íon
    \item Seleção de precursores
\end{itemize}

\column{0.5\textwidth}
\textbf{MS2 (Tandem MS):}
\begin{itemize}
    \item Fragmentação (CID, HCD, ETD)
    \item Espectro de fragmentos
    \item Sequência do peptídeo
\end{itemize}
\end{columns}
\end{frame}

\begin{frame}{Técnicas de Ionização}
\begin{columns}
\column{0.5\textwidth}
\begin{block}{ESI}
\textbf{Electrospray Ionization}
\begin{itemize}
    \item Ionização suave
    \item Íons múltiplos carregados
    \item Acoplado a LC (LC-MS)
    \item Peptídeos e proteínas pequenas
    \item Mais usado em proteômica
\end{itemize}
\end{block}

\column{0.5\textwidth}
\begin{block}{MALDI-TOF}
\textbf{Matrix-Assisted Laser Desorption/Ionization}
\begin{itemize}
    \item Laser dessorve amostra
    \item Matriz absorve energia
    \item Íons simples carregados
    \item Time-of-Flight (ToF)
    \item Bom para proteínas grandes
\end{itemize}
\end{block}
\end{columns}

\vspace{0.3cm}

\begin{block}{Analisadores de Massa}
\begin{itemize}
    \item \textbf{Orbitrap:} alta resolução, alta acurácia
    \item \textbf{Q-ToF:} tandem MS, bom para PTMs
    \item \textbf{Triple Quad:} quantificação (SRM/MRM)
\end{itemize}
\end{block}
\end{frame}

\begin{frame}{Quantificação Proteômica}
\begin{columns}
\column{0.5\textwidth}
\begin{block}{Label-Free}
\textbf{Sem Marcação}
\begin{itemize}
    \item Intensidade do sinal MS1
    \item Contagem de espectros (SpC)
    \item iBAQ, LFQ
    \item Simples, econômico
    \item Variabilidade maior
\end{itemize}
\end{block}

\begin{block}{SILAC}
\textbf{Stable Isotope Labeling}
\begin{itemize}
    \item Aminoácidos marcados ($^{13}$C, $^{15}$N)
    \item Incorporação metabólica
    \item Multiplexação (até 3 condições)
    \item Quantificação acurada
    \item Requer cultura celular
\end{itemize}
\end{block}

\column{0.5\textwidth}
\begin{block}{Marcação Química}
\textbf{TMT/iTRAQ}
\begin{itemize}
    \item Tags isobáricas
    \item Multiplexação (6-18 condições)
    \item Quantificação em MS2
    \item Versátil (qualquer amostra)
    \item Ratio compression
\end{itemize}
\end{block}

\vspace{0.3cm}

\begin{exampleblock}{Escolha do Método}
\begin{itemize}
    \item \textbf{Label-free:} descoberta, custo
    \item \textbf{SILAC:} acurácia, cultura
    \item \textbf{TMT:} multiplexação, clínica
\end{itemize}
\end{exampleblock}
\end{columns}
\end{frame}

\begin{frame}{Interações Proteína-Proteína}
\begin{block}{Métodos Experimentais}
\begin{columns}
\column{0.5\textwidth}
\textbf{Y2H (Yeast Two-Hybrid):}
\begin{itemize}
    \item In vivo em levedura
    \item Interações binárias
    \item Triagem em larga escala
    \item Falsos positivos
\end{itemize}

\vspace{0.3cm}

\textbf{Co-IP (Co-Imunoprecipitação):}
\begin{itemize}
    \item Anticorpo puxa proteína
    \item Complexos co-purificados
    \item Condições nativas
    \item Validação de interações
\end{itemize}

\column{0.5\textwidth}
\textbf{AP-MS (Affinity Purification-MS):}
\begin{itemize}
    \item Tag (FLAG, HA, Strep)
    \item Purificação de complexos
    \item Identificação por MS
    \item Mapeamento de complexos
\end{itemize}

\vspace{0.3cm}

\textbf{BioID/APEX:}
\begin{itemize}
    \item Marcação por proximidade
    \item Biotina em vizinhos
    \item Interações transitórias
    \item In vivo
\end{itemize}
\end{columns}
\end{block}
\end{frame}

\begin{frame}{Bancos de Dados de Interações}
\begin{columns}
\column{0.5\textwidth}
\begin{block}{STRING}
\url{https://string-db.org}
\begin{itemize}
    \item Interações conhecidas e preditas
    \item Integração de evidências
    \item Scores de confiança
    \item Visualização de redes
    \item Análise de enriquecimento
\end{itemize}
\end{block}

\begin{block}{BioGRID}
\begin{itemize}
    \item Interações experimentais
    \item Curado manualmente
    \item Múltiplas espécies
\end{itemize}
\end{block}

\column{0.5\textwidth}
\begin{block}{IntAct}
\begin{itemize}
    \item Base do EBI
    \item Dados estruturados
    \item PSI-MI format
\end{itemize}
\end{block}

\begin{block}{CORUM}
\begin{itemize}
    \item Complexos proteicos
    \item Mamíferos
    \item Composição e função
\end{itemize}
\end{block}

\vspace{0.3cm}

\importante{
Dados de diferentes métodos são complementares: integração é chave!
}
\end{columns}
\end{frame}

\begin{frame}[fragile]{Análise de Dados Proteômicos}
\begin{lstlisting}[language=Python, caption=Processamento de dados MS, basicstyle=\ttfamily\footnotesize]
import pandas as pd
import numpy as np
from scipy.stats import ttest_ind
import matplotlib.pyplot as plt

# Carregar dados de quantificacao (ex: MaxQuant output)
df = pd.read_csv('proteinGroups.txt', sep='\t')

# Selecionar colunas de intensidade
control_cols = [f'LFQ intensity Ctrl_{i}' for i in range(1,4)]
treatment_cols = [f'LFQ intensity Treat_{i}' for i in range(1,4)]

# Log2 transform
df[control_cols] = np.log2(df[control_cols].replace(0, np.nan))
df[treatment_cols] = np.log2(df[treatment_cols].replace(0, np.nan))

# Filtrar proteinas com muitos missing values
df = df.dropna(thresh=4, subset=control_cols+treatment_cols)

# Calcular fold change
df['log2FC'] = df[treatment_cols].mean(axis=1) - \
               df[control_cols].mean(axis=1)

# Teste t
t_stat, p_val = [], []
for idx, row in df.iterrows():
    ctrl = row[control_cols].dropna()
    treat = row[treatment_cols].dropna()
    if len(ctrl) >= 2 and len(treat) >= 2:
        t, p = ttest_ind(ctrl, treat)
        t_stat.append(t)
        p_val.append(p)
    else:
        t_stat.append(np.nan)
        p_val.append(np.nan)

df['p_value'] = p_val
df['-log10(p)'] = -np.log10(df['p_value'])
\end{lstlisting}
\end{frame}

\begin{frame}[fragile]{Visualização: Volcano Plot}
\begin{lstlisting}[language=Python, caption=Volcano plot proteômico, basicstyle=\ttfamily\footnotesize]
# Continuacao do codigo anterior

# Thresholds
fc_thresh = 1  # log2 fold change
p_thresh = 0.05

# Classificar proteinas
df['significant'] = 'No'
df.loc[(df['log2FC'] > fc_thresh) & (df['p_value'] < p_thresh),
       'significant'] = 'Up'
df.loc[(df['log2FC'] < -fc_thresh) & (df['p_value'] < p_thresh),
       'significant'] = 'Down'

# Volcano plot
plt.figure(figsize=(10, 8))
colors = {'No': 'gray', 'Up': 'red', 'Down': 'blue'}
for sig, group in df.groupby('significant'):
    plt.scatter(group['log2FC'], group['-log10(p)'],
                c=colors[sig], label=sig, alpha=0.6, s=20)

# Threshold lines
plt.axhline(y=-np.log10(p_thresh), color='black',
            linestyle='--', linewidth=1)
plt.axvline(x=fc_thresh, color='black',
            linestyle='--', linewidth=1)
plt.axvline(x=-fc_thresh, color='black',
            linestyle='--', linewidth=1)

plt.xlabel('log2(Fold Change)', fontsize=14)
plt.ylabel('-log10(p-value)', fontsize=14)
plt.title('Volcano Plot: Proteinas Diferencialmente Expressas')
plt.legend()
plt.grid(True, alpha=0.3)
\end{lstlisting}
\end{frame}

% ============================================================================
% SEÇÃO 7: EXERCÍCIOS
% ============================================================================
\section{Exercícios}

\begin{frame}{Exercícios - Parte 1: Cinética Enzimática}
\begin{block}{Questão 1: Parâmetros de MM}
Uma enzima apresenta $K_M = 5$ mM e $V_{max} = 100$ $\mu$mol/min. Calcule:
\begin{enumerate}
    \item A velocidade quando $[S] = 15$ mM
    \item A concentração de substrato para $v = 75$ $\mu$mol/min
    \item A eficiência catalítica se $[E]_0 = 1$ $\mu$M
\end{enumerate}
\end{block}

\begin{block}{Questão 2: Inibição}
Um inibidor competitivo tem $K_I = 2$ mM. Se $K_M = 5$ mM e $[I] = 10$ mM:
\begin{enumerate}
    \item Qual é o $K_M^{app}$?
    \item Como isso afeta o gráfico de Lineweaver-Burk?
    \item Que concentração de substrato é necessária para superar 90\% da inibição?
\end{enumerate}
\end{block}
\end{frame}

\begin{frame}{Exercícios - Parte 2: Estrutura e Análise}
\begin{block}{Questão 3: Ramachandran}
Explique por que:
\begin{enumerate}
    \item Glicina tem maior liberdade conformacional
    \item Prolina restringe ângulos $\phi$ e $\psi$
    \item Regiões proibidas existem no diagrama
\end{enumerate}
\end{block}

\begin{alertblock}{Questão 4: Projeto Prático - BioPython}
Escolha uma proteína do PDB:
\begin{enumerate}
    \item Baixe e parse a estrutura com BioPython
    \item Calcule proporção de estruturas secundárias ($\alpha$, $\beta$, loop)
    \item Identifique resíduos hidrofóbicos vs. hidrofílicos
    \item Visualize em PyMOL colorindo por propriedade
    \item Analise cavidades potenciais de ligação
\end{enumerate}
\end{alertblock}
\end{frame}

\begin{frame}{Exercícios - Parte 3: Sinalização e PTMs}
\begin{block}{Questão 5: Cascata MAPK}
Considerando amplificação 10x em cada nível:
\begin{enumerate}
    \item Quantas moléculas de ERK são ativadas por 1 RTK?
    \item Se 100 RTKs são ativados, quantas ERK fosforiladas?
    \item Como feedback negativo afeta isso?
\end{enumerate}
\end{block}

\begin{block}{Questão 6: PTMs}
Uma proteína tem os seguintes sítios:
\begin{itemize}
    \item Lys$_{48}$: pode ser acetilada ou ubiquitinada
    \item Ser$_{132}$: pode ser fosforilada ou O-GlcNAcilada
\end{itemize}
\begin{enumerate}
    \item Quantas combinações de PTMs são possíveis?
    \item Como crosstalk entre PTMs regula função?
    \item Que técnica usaria para detectar cada modificação?
\end{enumerate}
\end{block}
\end{frame}

\begin{frame}{Exercícios - Parte 4: Proteômica}
\begin{alertblock}{Projeto Final: Pipeline Proteômico}
Desenvolva análise completa de dados de proteômica:
\begin{enumerate}
    \item Baixe dados públicos (PRIDE, ProteomeXchange)
    \item Pré-processamento: log-transform, normalização, filtro
    \item Análise estatística: teste t, correção FDR
    \item Volcano plot e heatmap
    \item Análise de enriquecimento GO/KEGG
    \item Rede de interações proteína-proteína (STRING)
    \item Identificação de hubs e módulos funcionais
    \item Relatório com interpretação biológica
\end{enumerate}
\end{alertblock}

\begin{block}{Questão 7: Quantificação}
Compare métodos label-free, SILAC e TMT:
\begin{itemize}
    \item Quando usar cada um?
    \item Vantagens e limitações?
    \item Qual é mais apropriado para amostras clínicas?
\end{itemize}
\end{block}
\end{frame}

% ============================================================================
% SEÇÃO 8: REFERÊNCIAS
% ============================================================================
\section{Referências}

\begin{frame}{Referências Principais}
\begin{thebibliography}{99}
\bibitem{berg} Berg J.M., Tymoczko J.L., Stryer L. (2019). \textit{Biochemistry}, 9th ed. W.H. Freeman.

\bibitem{voet} Voet D., Voet J.G., Pratt C.W. (2016). \textit{Fundamentals of Biochemistry}, 5th ed. Wiley.

\bibitem{lodish} Lodish H. et al. (2021). \textit{Molecular Cell Biology}, 9th ed. W.H. Freeman.

\bibitem{aebersold} Aebersold R., Mann M. (2016). Mass-spectrometric exploration of proteome structure and function. \textit{Nature}.

\bibitem{mann} Mann M., Jensen O.N. (2003). Proteomic analysis of post-translational modifications. \textit{Nature Biotechnology}.
\end{thebibliography}
\end{frame}

\begin{frame}{Referências Complementares}
\begin{thebibliography}{99}
\bibitem{cornish} Cornish-Bowden A. (2012). \textit{Fundamentals of Enzyme Kinetics}, 4th ed. Wiley-Blackwell.

\bibitem{cox} Cox J., Mann M. (2011). Quantitative, high-resolution proteomics for data-driven systems biology. \textit{Annual Review of Biochemistry}.

\bibitem{koshland} Koshland D.E. et al. (1966). Comparison of experimental binding data and theoretical models in proteins containing subunits. \textit{Biochemistry}.

\bibitem{goldbeter} Goldbeter A., Koshland D.E. (1981). An amplified sensitivity arising from covalent modification in biological systems. \textit{PNAS}.

\bibitem{ptm} Venne A.S. et al. (2014). The next level of complexity: Crosstalk of posttranslational modifications. \textit{Proteomics}.
\end{thebibliography}
\end{frame}

\begin{frame}{Leitura Complementar e Recursos}
\begin{block}{Livros Especializados}
\begin{itemize}
    \item Bairoch A. et al. (2016). \textit{The SWISS-PROT Protein Knowledgebase}
    \item Liebler D.C. (2001). \textit{Introduction to Proteomics}
    \item Cech T.R., Atkins J.F. (2011). \textit{RNA Worlds}, 3rd ed.
\end{itemize}
\end{block}

\begin{block}{Recursos Online}
\begin{itemize}
    \item \textbf{PDB:} \url{https://www.rcsb.org}
    \item \textbf{UniProt:} \url{https://www.uniprot.org}
    \item \textbf{PhosphoSitePlus:} \url{https://www.phosphosite.org}
    \item \textbf{STRING:} \url{https://string-db.org}
    \item \textbf{PRIDE:} \url{https://www.ebi.ac.uk/pride}
    \item \textbf{ProteomeXchange:} \url{http://www.proteomexchange.org}
    \item \textbf{BRENDA:} \url{https://www.brenda-enzymes.org} (enzimas)
\end{itemize}
\end{block}
\end{frame}

% ============================================================================
% SLIDE FINAL
% ============================================================================
\begin{frame}
\begin{center}
{\Huge Obrigado!}

\vspace{1cm}

{\Large Capítulo 7: Protein Systems}

\vspace{1cm}

\textbf{Próximo Capítulo:}\\
Metabolic Systems (Sistemas Metabólicos)

\vspace{1cm}

\textit{Dúvidas? Perguntas?}
\end{center}
\end{frame}

\end{document}
